\documentclass[11pt,a4paper]{article}

%====================================================
% PACKAGES
%====================================================
\usepackage[utf8]{inputenc}
\usepackage[T1]{fontenc}
\usepackage[french]{babel}
\usepackage[a4paper,margin=2cm]{geometry}
\usepackage{lmodern}
\usepackage{microtype}
\usepackage{amsmath,amssymb,amsthm,mathtools}
\usepackage{enumitem}
\usepackage{fancyhdr}
\usepackage{lastpage}
\usepackage{xcolor}
\usepackage[most]{tcolorbox}
\usepackage{hyperref}

%====================================================
% COULEURS ET STYLE
%====================================================
\definecolor{bleuprepa}{RGB}{20,55,110}
\definecolor{rougeprepa}{RGB}{150,35,35}
\definecolor{vertprepa}{RGB}{25,115,75}
\definecolor{orangeprepa}{RGB}{190,110,20}
\definecolor{grisclair}{RGB}{245,247,250}

\hypersetup{
  colorlinks=true,
  linkcolor=bleuprepa,
  urlcolor=bleuprepa
}

\pagestyle{fancy}
\fancyhf{}
\fancyhead[L]{\textsc{CPGE scientifiques}}
\fancyhead[C]{\textsc{TD -- Suites numériques}}
\fancyhead[R]{\textsc{MPSI/PCSI/MP}}
\fancyfoot[C]{Page \thepage/\pageref{LastPage}}

\setlength{\parindent}{0pt}
\setlength{\parskip}{0.4em}

\setlist[itemize]{label=\(\triangleright\),leftmargin=1.2cm}
\setlist[enumerate]{label=\textbf{\arabic*.},leftmargin=1.1cm}

\tcbset{
  enhanced,
  breakable,
  sharp corners,
  boxrule=0.7pt,
  colback=white,
  fonttitle=\bfseries
}

\newcounter{exo}
\newcounter{sol}

\newcommand{\etoiles}[1]{\hfill{\small\textcolor{orangeprepa}{#1}}}

\newtcolorbox{exercicebox}[2]{
  colframe=bleuprepa,
  colback=blue!2!white,
  title={Exercice #1 -- #2}
}

\newtcolorbox{solutionbox}[2]{
  colframe=vertprepa,
  colback=green!2!white,
  title={Correction #1 -- #2}
}

\newcommand{\exercice}[2]{
  \refstepcounter{exo}
  \begin{exercicebox}{\theexo}{#1}
  #2
  \end{exercicebox}
}

\newcommand{\solution}[2]{
  \refstepcounter{sol}
  \begin{solutionbox}{\thesol}{#1}
  #2
  \end{solutionbox}
}

\newcommand{\N}{\mathbb{N}}
\newcommand{\R}{\mathbb{R}}
\newcommand{\C}{\mathbb{C}}
\newcommand{\eps}{\varepsilon}

%====================================================
% DOCUMENT
%====================================================
\begin{document}

\begin{center}
{\Huge \bfseries TD complet -- Suites numériques}\\[0.2cm]
{\Large Convergence, monotonie, suites extraites, récurrences et comparaisons}\\[0.2cm]
{\large Classes préparatoires scientifiques -- MPSI, PCSI, MP}\\[0.3cm]
\textbf{Document prêt à compiler -- Énoncés et corrigés détaillés}
\end{center}

\vspace{0.4cm}

\begin{tcolorbox}[colframe=rougeprepa,colback=red!2!white,title={Objectifs du TD}]
Ce TD vise à maîtriser les méthodes fondamentales sur les suites numériques :
\begin{itemize}
  \item démontrer une convergence à partir de la définition ;
  \item calculer des limites par opérations, encadrement, équivalents ;
  \item utiliser monotonie, bornes et théorème de convergence monotone ;
  \item étudier des suites récurrentes \(u_{n+1}=f(u_n)\) ;
  \item exploiter les suites extraites pour prouver convergence ou divergence ;
  \item traiter des problèmes longs de niveau concours.
\end{itemize}
\end{tcolorbox}

\tableofcontents
\newpage

%====================================================
\section{Exercices d'application directe -- APP}

\exercice{Définition d'une suite et premières limites \etoiles{★}}{
On considère les suites définies, pour \(n\geq 1\), par
\[
u_n=\frac{1}{n},\qquad v_n=\frac{n}{n+1},\qquad w_n=(-1)^n.
\]
\begin{enumerate}
  \item Donner explicitement les cinq premiers termes de chaque suite.
  \item Montrer que \((u_n)\) converge vers \(0\) en utilisant la définition avec \(\eps\).
  \item Montrer que \((v_n)\) converge vers \(1\).
  \item Montrer que \((w_n)\) ne converge pas.
\end{enumerate}
}

\exercice{Suites bornées \etoiles{★}}{
Pour \(n\geq 1\), on pose
\[
u_n=\frac{2n+1}{n+3}.
\]
\begin{enumerate}
  \item Montrer que \(u_n=2-\dfrac{5}{n+3}\).
  \item En déduire un encadrement simple de \(u_n\).
  \item Montrer que \((u_n)\) est bornée.
  \item Déterminer sa limite.
\end{enumerate}
}

\exercice{Limite par division par le terme dominant \etoiles{★}}{
Déterminer les limites suivantes :
\[
u_n=\frac{3n^2-2n+1}{5n^2+n-7},\qquad
v_n=\frac{4n^3+n}{2n^4+1},\qquad
w_n=\frac{-n^5+3n^2}{2n^5+n}.
\]
Pour chaque suite, on justifiera soigneusement le passage à la limite.
}

\exercice{Utilisation du théorème des gendarmes \etoiles{★}}{
Pour \(n\geq 1\), on pose
\[
u_n=\frac{\sin n}{n},\qquad v_n=\frac{\cos(n^2)}{\sqrt n}.
\]
\begin{enumerate}
  \item Encadrer \(u_n\) entre deux suites simples.
  \item En déduire la limite de \(u_n\).
  \item Encadrer \(v_n\) entre deux suites simples.
  \item En déduire la limite de \(v_n\).
\end{enumerate}
}

\exercice{Radicaux et quantité conjuguée \etoiles{★★}}{
Déterminer les limites suivantes :
\[
u_n=\sqrt{n^2+3n}-n,\qquad
v_n=\sqrt{4n^2+n}-2n,\qquad
w_n=\sqrt{n^2+1}-\sqrt{n^2-1}.
\]
On utilisera une quantité conjuguée.
}

\exercice{Suites géométriques \etoiles{★}}{
Soit \(q\in\R\). On considère la suite \(u_n=q^n\).
\begin{enumerate}
  \item Étudier la limite de \((u_n)\) lorsque \(|q|<1\).
  \item Étudier le cas \(q=1\).
  \item Étudier le cas \(q=-1\).
  \item Étudier le cas \(q>1\).
  \item Étudier le cas \(q<-1\).
\end{enumerate}
}

\exercice{Passage à la limite dans les inégalités \etoiles{★★}}{
Soient \((u_n)\) et \((v_n)\) deux suites réelles telles que
\[
\forall n\in\N,\qquad u_n\leq v_n.
\]
On suppose que \(u_n\to \ell\) et \(v_n\to m\).
\begin{enumerate}
  \item Montrer que \(\ell\leq m\).
  \item Donner un exemple où \(u_n<v_n\) pour tout \(n\), mais où \(\ell=m\).
  \item Expliquer pourquoi on ne peut pas conclure à une inégalité stricte sur les limites.
\end{enumerate}
}

\exercice{Convergence implique bornitude \etoiles{★★}}{
Soit \((u_n)\) une suite réelle convergeant vers \(\ell\in\R\).
\begin{enumerate}
  \item En prenant \(\eps=1\), montrer qu'il existe un rang \(N\) tel que \(|u_n-\ell|\leq 1\) pour \(n\geq N\).
  \item En déduire que la suite est bornée à partir du rang \(N\).
  \item Conclure que \((u_n)\) est bornée sur \(\N\).
  \item La réciproque est-elle vraie ? Justifier.
\end{enumerate}
}

\exercice{Monotonie par différence \etoiles{★}}{
Étudier la monotonie des suites suivantes :
\[
u_n=\frac{n}{n+1},\qquad v_n=n^2-3n,\qquad w_n=\frac{1}{n+1}.
\]
Pour chacune, on calculera \(u_{n+1}-u_n\), \(v_{n+1}-v_n\) ou \(w_{n+1}-w_n\).
}

\exercice{Monotonie par quotient \etoiles{★★}}{
Pour \(n\geq 1\), on pose
\[
u_n=\frac{2^n}{n!}.
\]
\begin{enumerate}
  \item Calculer \(\dfrac{u_{n+1}}{u_n}\).
  \item Montrer que \((u_n)\) est croissante jusqu'à un certain rang puis décroissante.
  \item Montrer que la suite est bornée.
  \item Déterminer sa limite.
\end{enumerate}
}

%====================================================
\section{Exercices de méthode -- METH}

\exercice{Méthode \(\eps\) \etoiles{★★}}{
Montrer directement avec la définition que
\[
\frac{2n+1}{3n-4}\longrightarrow \frac23.
\]
\begin{enumerate}
  \item Calculer \(\left|\dfrac{2n+1}{3n-4}-\dfrac23\right|\).
  \item Trouver une condition suffisante sur \(n\) pour que cette quantité soit inférieure à \(\eps\).
  \item Conclure rigoureusement.
\end{enumerate}
}

\exercice{Limites infinies \etoiles{★★}}{
Montrer avec la définition que
\[
n^2-3n\longrightarrow +\infty.
\]
\begin{enumerate}
  \item Rappeler la définition de \(u_n\to+\infty\).
  \item Montrer que, pour \(n\geq 6\), on a \(n^2-3n\geq \dfrac{n^2}{2}\).
  \item Conclure.
\end{enumerate}
}

\exercice{Somme, produit, quotient \etoiles{★★}}{
On suppose que
\[
u_n\to 2,\qquad v_n\to -3.
\]
Déterminer les limites de
\[
a_n=3u_n-2v_n,\qquad b_n=u_nv_n,\qquad c_n=\frac{u_n+1}{v_n-1}.
\]
Justifier que le quotient est défini à partir d'un certain rang.
}

\exercice{Limite et valeur absolue \etoiles{★★}}{
Soit \((u_n)\) une suite réelle.
\begin{enumerate}
  \item Montrer que si \(u_n\to 0\), alors \(|u_n|\to 0\).
  \item Montrer que si \(|u_n|\to 0\), alors \(u_n\to 0\).
  \item La propriété reste-t-elle vraie pour une limite \(\ell\neq 0\) sous la forme \(|u_n|\to |\ell|\) ?
\end{enumerate}
}

\exercice{Suite encadrée par deux suites convergentes \etoiles{★★}}{
Soit \((u_n)\) une suite telle que
\[
\frac{n}{n+1}\leq u_n\leq \frac{n+2}{n+1}.
\]
\begin{enumerate}
  \item Déterminer les limites des deux suites encadrantes.
  \item En déduire la limite de \((u_n)\).
  \item Donner un exemple explicite d'une telle suite \(u_n\).
\end{enumerate}
}

\exercice{Étude complète d'une suite explicite \etoiles{★★}}{
Pour \(n\geq 1\), on pose
\[
u_n=\frac{n^2+1}{n^2+n+1}.
\]
\begin{enumerate}
  \item Montrer que \(0<u_n<1\).
  \item Étudier la monotonie de \((u_n)\).
  \item Montrer que \((u_n)\) converge.
  \item Déterminer sa limite.
\end{enumerate}
}

\exercice{Suite alternée simple \etoiles{★★}}{
On pose
\[
u_n=(-1)^n+\frac1n,\qquad n\geq 1.
\]
\begin{enumerate}
  \item Étudier les deux suites extraites \((u_{2n})\) et \((u_{2n+1})\).
  \item Déterminer leurs limites respectives.
  \item En déduire que \((u_n)\) diverge.
\end{enumerate}
}

\exercice{Recouvrement pair-impair \etoiles{★★}}{
Soit \((u_n)\) une suite réelle. On suppose que
\[
u_{2n}\to \ell,\qquad u_{2n+1}\to \ell.
\]
\begin{enumerate}
  \item Rédiger une preuve complète de \(u_n\to \ell\).
  \item Appliquer ce résultat à la suite définie par
  \[
  u_{2n}=\frac{1}{n+1},\qquad u_{2n+1}=\frac{2}{n+1}.
  \]
\end{enumerate}
}

\exercice{Bornes supérieure et inférieure \etoiles{★★★}}{
On considère
\[
A=\left\{1-\frac1n\;:\; n\in\N^*\right\}.
\]
\begin{enumerate}
  \item Montrer que \(1\) est un majorant de \(A\).
  \item Construire une suite d'éléments de \(A\) convergeant vers \(1\).
  \item En déduire que \(\sup A=1\).
  \item Déterminer \(\inf A\).
\end{enumerate}
}

\exercice{Suite définie par une somme \etoiles{★★★}}{
Pour \(n\geq 1\), on pose
\[
u_n=\sum_{k=1}^n \frac1{k(k+1)}.
\]
\begin{enumerate}
  \item Décomposer \(\dfrac1{k(k+1)}\) en éléments simples.
  \item Simplifier la somme.
  \item Étudier la monotonie de \((u_n)\).
  \item Déterminer sa limite.
\end{enumerate}
}

%====================================================
\section{Exercices de démonstration -- DEM}

\exercice{Unicité de la limite \etoiles{★★}}{
Soit \((u_n)\) une suite réelle.
\begin{enumerate}
  \item Supposer que \(u_n\to \ell\) et \(u_n\to m\).
  \item Montrer que pour tout \(\eps>0\), on a \(|\ell-m|\leq \eps\).
  \item Conclure que \(\ell=m\).
\end{enumerate}
}

\exercice{Théorème des gendarmes \etoiles{★★★}}{
Soient \((a_n)\), \((u_n)\), \((b_n)\) trois suites réelles telles que, à partir d'un certain rang,
\[
a_n\leq u_n\leq b_n.
\]
On suppose que \(a_n\to \ell\) et \(b_n\to \ell\).
\begin{enumerate}
  \item Rappeler la définition de \(a_n\to\ell\) et \(b_n\to\ell\).
  \item Montrer que, pour \(n\) assez grand, \(\ell-\eps\leq u_n\leq \ell+\eps\).
  \item Conclure.
\end{enumerate}
}

\exercice{Théorème de convergence monotone \etoiles{★★★}}{
Soit \((u_n)\) une suite réelle croissante et majorée.
\begin{enumerate}
  \item Poser \(A=\{u_n:n\in\N\}\). Justifier que \(A\) admet une borne supérieure.
  \item Noter \(\ell=\sup A\). Montrer que \(u_n\leq \ell\) pour tout \(n\).
  \item Montrer que pour tout \(\eps>0\), il existe \(N\) tel que \(\ell-\eps<u_N\leq \ell\).
  \item Utiliser la croissance pour montrer que \(u_n\to \ell\).
\end{enumerate}
}

\exercice{Version divergente du théorème monotone \etoiles{★★★}}{
Soit \((u_n)\) une suite croissante non majorée.
\begin{enumerate}
  \item Rappeler la définition de \(u_n\to+\infty\).
  \item Soit \(A\in\R\). Utiliser le fait que \((u_n)\) n'est pas majorée.
  \item Conclure que \(u_n\to+\infty\).
\end{enumerate}
}

\exercice{Suites adjacentes \etoiles{★★★}}{
Soient \((u_n)\) et \((v_n)\) deux suites réelles telles que :
\[
(u_n)\text{ est croissante},\qquad (v_n)\text{ est décroissante},\qquad u_n\leq v_n,
\]
et
\[
v_n-u_n\to 0.
\]
\begin{enumerate}
  \item Montrer que \((u_n)\) est majorée.
  \item Montrer que \((v_n)\) est minorée.
  \item En déduire que les deux suites convergent.
  \item Montrer qu'elles ont la même limite.
\end{enumerate}
}

\exercice{Suites extraites d'une suite convergente \etoiles{★★}}{
Soit \((u_n)\) une suite convergeant vers \(\ell\), et soit \(\varphi:\N\to\N\) strictement croissante.
\begin{enumerate}
  \item Montrer par récurrence que \(\varphi(n)\geq n\).
  \item En déduire que \(u_{\varphi(n)}\to \ell\).
\end{enumerate}
}

\exercice{Critère de divergence par extraction \etoiles{★★}}{
Soit \((u_n)\) une suite réelle.
\begin{enumerate}
  \item Montrer que si deux suites extraites de \((u_n)\) convergent vers deux limites distinctes, alors \((u_n)\) diverge.
  \item Appliquer ce résultat à \(u_n=(-1)^n+\dfrac{1}{n+1}\).
\end{enumerate}
}

\exercice{Bolzano-Weierstrass : usage \etoiles{★★★}}{
Soit \((u_n)\) une suite réelle bornée.
\begin{enumerate}
  \item Énoncer le théorème de Bolzano-Weierstrass.
  \item Montrer que \((u_n)\) possède au moins une suite extraite convergente.
  \item Donner un exemple montrant que \((u_n)\) n'est pas nécessairement convergente.
\end{enumerate}
}

\exercice{Caractérisation séquentielle du supremum \etoiles{★★★}}{
Soit \(A\subset\R\) non vide et majorée, et soit \(s\in\R\).
\begin{enumerate}
  \item Montrer que si \(s=\sup A\), alors il existe une suite \((a_n)\) d'éléments de \(A\) telle que \(a_n\to s\).
  \item Réciproquement, montrer que si \(s\) majore \(A\) et s'il existe une suite \((a_n)\) d'éléments de \(A\) telle que \(a_n\to s\), alors \(s=\sup A\).
\end{enumerate}
}

\exercice{Comparaison asymptotique : définitions \etoiles{★★}}{
Soient \((u_n)\) et \((v_n)\) deux suites réelles avec \(v_n\neq0\) à partir d'un certain rang.
\begin{enumerate}
  \item Donner la définition de \(u_n=o(v_n)\).
  \item Donner la définition de \(u_n=O(v_n)\).
  \item Donner la définition de \(u_n\sim v_n\).
  \item Montrer que \(u_n\sim v_n\) implique \(u_n=O(v_n)\).
\end{enumerate}
}

%====================================================
\section{Exercices de réflexion -- PREP}

\exercice{Une suite bornée dont les extraites convergent différemment \etoiles{★★★}}{
On pose
\[
u_n=\sin\left(\frac{n\pi}{2}\right).
\]
\begin{enumerate}
  \item Calculer \(u_{4n}\), \(u_{4n+1}\), \(u_{4n+2}\), \(u_{4n+3}\).
  \item Montrer que \((u_n)\) est bornée.
  \item Montrer que \((u_n)\) diverge.
  \item Déterminer l'ensemble des valeurs prises par la suite.
\end{enumerate}
}

\exercice{Suite récurrente affine \etoiles{★★}}{
Soit \(a\in]-1,1[\) et \(b\in\R\). On définit
\[
u_{n+1}=au_n+b.
\]
\begin{enumerate}
  \item Résoudre l'équation du point fixe \(\ell=a\ell+b\).
  \item Poser \(v_n=u_n-\ell\). Montrer que \(v_{n+1}=av_n\).
  \item En déduire l'expression de \(u_n\).
  \item Déterminer la limite de \((u_n)\).
\end{enumerate}
}

\exercice{Suite récurrente contractante simple \etoiles{★★★}}{
Soit \((u_n)\) définie par
\[
u_0\in[0,1],\qquad u_{n+1}=\frac{u_n+1}{2}.
\]
\begin{enumerate}
  \item Montrer que \([0,1]\) est stable par \(f(x)=\dfrac{x+1}{2}\).
  \item Montrer que \((u_n)\) est croissante.
  \item Montrer que \((u_n)\) converge.
  \item Déterminer sa limite.
\end{enumerate}
}

\exercice{Suite de Héron pour \(\sqrt2\) \etoiles{★★★★}}{
Soit \(u_0>\sqrt2\), et
\[
u_{n+1}=\frac12\left(u_n+\frac2{u_n}\right).
\]
\begin{enumerate}
  \item Montrer que \(u_n>\sqrt2\) pour tout \(n\).
  \item Montrer que \((u_n)\) est décroissante.
  \item En déduire que \((u_n)\) converge.
  \item Déterminer sa limite.
\end{enumerate}
}

\exercice{Deux suites couplées \etoiles{★★★★}}{
Soient \(a>0\), \(b>0\), avec \(a<b\). On définit
\[
u_0=a,\qquad v_0=b,\qquad
u_{n+1}=\frac{2u_nv_n}{u_n+v_n},\qquad
v_{n+1}=\frac{u_n+v_n}{2}.
\]
\begin{enumerate}
  \item Montrer que \(0<u_n\leq v_n\) pour tout \(n\).
  \item Montrer que \((u_n)\) est croissante et que \((v_n)\) est décroissante.
  \item Montrer que \(v_{n+1}-u_{n+1}=\dfrac{(v_n-u_n)^2}{2(u_n+v_n)}\).
  \item En déduire que les deux suites sont adjacentes.
\end{enumerate}
}

\exercice{Suite définie par un maximum \etoiles{★★★}}{
Soit
\[
u_n=\max\left(\frac1n,\frac{(-1)^n}{2}\right),\qquad n\geq1.
\]
\begin{enumerate}
  \item Calculer \(u_{2n}\) et \(u_{2n+1}\).
  \item Déterminer leurs limites.
  \item La suite \((u_n)\) converge-t-elle ?
\end{enumerate}
}

\exercice{Suite et partie entière \etoiles{★★★}}{
Pour \(n\geq1\), on pose
\[
u_n=\frac{\lfloor \sqrt n\rfloor}{\sqrt n}.
\]
\begin{enumerate}
  \item Montrer que \(0<u_n\leq1\).
  \item Montrer que, pour tout \(n\), \(\lfloor\sqrt n\rfloor>\sqrt n-1\).
  \item En déduire un encadrement de \(u_n\).
  \item Déterminer la limite de \((u_n)\).
\end{enumerate}
}

\exercice{Comparaison polynomiale \etoiles{★★}}{
Déterminer un équivalent simple de chacune des suites suivantes :
\[
u_n=3n^4-2n^2+n,\qquad
v_n=\frac{5n^3-n+7}{2n^2+1},\qquad
w_n=\frac{n^2+1}{n^4+n}.
\]
En déduire leur limite.
}

\exercice{Équivalents et racines \etoiles{★★★}}{
Déterminer un équivalent simple de
\[
u_n=\sqrt{n^2+n+1},\qquad
v_n=\sqrt{n+1}-\sqrt n,\qquad
w_n=\sqrt{n^2+2n}-n.
\]
}

\exercice{Contre-exemple sur les équivalents \etoiles{★★★★}}{
On rappelle qu'on ne peut pas additionner des équivalents sans précaution.
\begin{enumerate}
  \item Donner deux suites \(u_n\sim v_n\) et \(a_n\sim b_n\).
  \item Choisir ces suites de sorte que \(u_n+a_n\) et \(v_n+b_n\) n'aient pas le même comportement asymptotique.
  \item Expliquer précisément où se situe la compensation.
\end{enumerate}
}

%====================================================
\section{Exercices type oraux et problèmes longs -- Concours}

\exercice{Oral CCP -- Suite définie par \(u_{n+1}=u_n-u_n^2\) \etoiles{★★★★}}{
Soit \(u_0\in]0,1[\), et
\[
u_{n+1}=u_n-u_n^2.
\]
\begin{enumerate}
  \item Montrer que \(0<u_n<1\) pour tout \(n\).
  \item Étudier la monotonie de \((u_n)\).
  \item Montrer que \((u_n)\) converge et déterminer sa limite.
  \item Montrer que
  \[
  \frac{1}{u_{n+1}}-\frac{1}{u_n}=\frac{1}{1-u_n}.
  \]
  \item En déduire que \(u_n\sim \dfrac1n\).
\end{enumerate}
}

\exercice{Oral Mines -- Suite logistique simple \etoiles{★★★★}}{
Soit \(u_0\in]0,1[\), et
\[
u_{n+1}=u_n(2-u_n).
\]
\begin{enumerate}
  \item Montrer que \(0<u_n<1\) pour tout \(n\).
  \item Montrer que \((u_n)\) est croissante.
  \item En déduire que \((u_n)\) converge.
  \item Déterminer sa limite.
  \item Étudier la suite \(v_n=1-u_n\).
\end{enumerate}
}

\exercice{Oral Centrale -- Moyennes arithmético-géométriques \etoiles{★★★★★}}{
Soient \(a,b>0\), \(a<b\). On définit
\[
u_0=a,\qquad v_0=b,\qquad
u_{n+1}=\sqrt{u_nv_n},\qquad v_{n+1}=\frac{u_n+v_n}{2}.
\]
\begin{enumerate}
  \item Montrer que \(0<u_n\leq v_n\) pour tout \(n\).
  \item Montrer que \((u_n)\) est croissante.
  \item Montrer que \((v_n)\) est décroissante.
  \item Montrer que \(v_n-u_n\to0\).
  \item Conclure que les deux suites convergent vers une même limite.
\end{enumerate}
}

\exercice{Problème long -- Construction du nombre \(e\) \etoiles{★★★★★}}{
On pose, pour \(n\geq1\),
\[
u_n=\left(1+\frac1n\right)^n,\qquad
v_n=\left(1+\frac1n\right)^{n+1}.
\]
On admet l'inégalité de Bernoulli : pour \(x>-1\) et \(m\in\N\), \((1+x)^m\geq1+mx\).
\begin{enumerate}
  \item Montrer que \((u_n)\) est croissante.
  \item Montrer que \((v_n)\) est décroissante.
  \item Montrer que \(u_n\leq v_n\).
  \item Montrer que \(v_n-u_n=\dfrac{u_n}{n}\).
  \item Montrer que \((u_n)\) est majorée.
  \item En déduire que \((u_n)\) et \((v_n)\) sont adjacentes.
\end{enumerate}
}

\exercice{Problème long -- Récurrence et point fixe \etoiles{★★★★}}{
Soit \(f(x)=\sqrt{2+x}\), et soit \((u_n)\) définie par
\[
u_0\in[0,2],\qquad u_{n+1}=f(u_n).
\]
\begin{enumerate}
  \item Montrer que \([0,2]\) est stable par \(f\).
  \item Montrer que \(u_{n+1}-u_n\) a le signe de \(2-u_n\).
  \item En déduire la monotonie de \((u_n)\).
  \item Montrer que \((u_n)\) converge.
  \item Déterminer sa limite.
\end{enumerate}
}

\exercice{Problème long -- Suites extraites et valeurs d'adhérence \etoiles{★★★★}}{
Soit \((u_n)\) une suite réelle bornée.
On suppose que les deux suites extraites \((u_{2n})\) et \((u_{2n+1})\) convergent.
\begin{enumerate}
  \item Montrer que si elles ont la même limite, alors \((u_n)\) converge.
  \item Montrer que si elles ont des limites distinctes, alors \((u_n)\) diverge.
  \item Appliquer au cas \(u_n=(-1)^n+\dfrac{1}{n}\).
  \item Appliquer au cas \(u_n=\dfrac{1+(-1)^n}{n+1}\).
\end{enumerate}
}

\exercice{Problème long -- Une suite implicite par encadrement \etoiles{★★★★}}{
Pour \(n\geq1\), on pose
\[
u_n=\sum_{k=1}^n \frac{1}{n+k}.
\]
\begin{enumerate}
  \item Montrer que
  \[
  \frac{n}{2n}\leq u_n\leq \frac{n}{n+1}.
  \]
  \item Ce premier encadrement suffit-il à déterminer la limite ?
  \item Montrer que, pour tout \(k\in\{1,\dots,n\}\),
  \[
  \frac1{2n}\leq \frac1{n+k}\leq \frac1n.
  \]
  \item Obtenir un meilleur encadrement à l'aide d'une comparaison plus fine :
  \[
  \frac{n}{2n}\leq u_n\leq \frac{n}{n+1}.
  \]
  \item Commenter la difficulté et expliquer pourquoi une méthode d'intégrales serait plus adaptée hors du cadre strict de ce TD.
\end{enumerate}
}

\exercice{Oral -- Suite complexe \etoiles{★★★}}{
Soit \(z\in\C\), et soit \(u_n=z^n\).
\begin{enumerate}
  \item Étudier le cas \(|z|<1\).
  \item Étudier le cas \(z=1\).
  \item Étudier le cas \(|z|>1\).
  \item Donner un exemple avec \(|z|=1\) où la suite diverge.
\end{enumerate}
}

\exercice{Oral -- Suite définie par \(u_{n+1}=\sqrt{u_n}\) \etoiles{★★★}}{
Soit \(u_0\in]0,+\infty[\), et
\[
u_{n+1}=\sqrt{u_n}.
\]
\begin{enumerate}
  \item Montrer que \(u_n>0\) pour tout \(n\).
  \item Étudier séparément les cas \(u_0>1\), \(u_0=1\), \(0<u_0<1\).
  \item Montrer que \((u_n)\) converge.
  \item Déterminer sa limite.
\end{enumerate}
}

\exercice{Oral -- Une suite rationnelle \etoiles{★★★★}}{
Soit
\[
u_n=\frac{1+2+\cdots+n}{n^2}.
\]
\begin{enumerate}
  \item Simplifier \(1+2+\cdots+n\).
  \item Déterminer une expression simple de \(u_n\).
  \item Étudier la monotonie de \((u_n)\).
  \item Déterminer sa limite.
\end{enumerate}
}

%====================================================
\newpage
\section{Corrigé complet}

\solution{Définition d'une suite et premières limites}{
Les premiers termes sont :
\[
u_1=1,\ u_2=\frac12,\ u_3=\frac13,\ u_4=\frac14,\ u_5=\frac15,
\]
\[
v_1=\frac12,\ v_2=\frac23,\ v_3=\frac34,\ v_4=\frac45,\ v_5=\frac56,
\]
et
\[
w_1=-1,\ w_2=1,\ w_3=-1,\ w_4=1,\ w_5=-1.
\]
Pour \(u_n\), soit \(\eps>0\). On choisit \(N\in\N^*\) tel que \(N\geq 1/\eps\). Alors, pour \(n\geq N\),
\[
|u_n-0|=\frac1n\leq \frac1N\leq \eps.
\]
Donc \(u_n\to0\).

Pour \(v_n\),
\[
\left|\frac{n}{n+1}-1\right|=\frac1{n+1}.
\]
Si \(n\geq N\geq 1/\eps\), alors \(\frac1{n+1}\leq\eps\). Donc \(v_n\to1\).

Enfin, \(w_{2n}=1\) et \(w_{2n+1}=-1\). Deux suites extraites convergent vers deux limites distinctes, donc \((w_n)\) diverge.
}

\solution{Suites bornées}{
On écrit
\[
\frac{2n+1}{n+3}=\frac{2(n+3)-5}{n+3}=2-\frac5{n+3}.
\]
Comme \(n\geq1\), on a \(n+3\geq4\), donc
\[
0<\frac5{n+3}\leq \frac54.
\]
Ainsi
\[
2-\frac54\leq u_n<2,
\]
c'est-à-dire
\[
\frac34\leq u_n<2.
\]
La suite est donc bornée. De plus,
\[
\frac5{n+3}\to0,
\]
donc
\[
u_n\to2.
\]
}

\solution{Limite par division par le terme dominant}{
On divise par la puissance dominante.

Pour \(u_n\),
\[
u_n=\frac{3-\frac2n+\frac1{n^2}}{5+\frac1n-\frac7{n^2}}\to \frac35.
\]

Pour \(v_n\),
\[
v_n=\frac{\frac4n+\frac1{n^3}}{2+\frac1{n^4}}\to0.
\]

Pour \(w_n\),
\[
w_n=\frac{-1+\frac3{n^3}}{2+\frac1{n^4}}\to -\frac12.
\]
Chaque passage est justifié par les opérations usuelles sur les limites et par le fait que les dénominateurs convergent vers des réels non nuls.
}

\solution{Utilisation du théorème des gendarmes}{
Pour tout \(n\geq1\),
\[
-1\leq \sin n\leq1.
\]
En divisant par \(n>0\),
\[
-\frac1n\leq \frac{\sin n}{n}\leq \frac1n.
\]
Les deux bornes tendent vers \(0\). Donc
\[
\frac{\sin n}{n}\to0.
\]

De même,
\[
-1\leq \cos(n^2)\leq1,
\]
donc
\[
-\frac1{\sqrt n}\leq \frac{\cos(n^2)}{\sqrt n}\leq \frac1{\sqrt n}.
\]
Les deux bornes tendent vers \(0\), donc
\[
v_n\to0.
\]
}

\solution{Radicaux et quantité conjuguée}{
On rationalise.

\[
u_n=\sqrt{n^2+3n}-n
=\frac{3n}{\sqrt{n^2+3n}+n}
=\frac3{\sqrt{1+\frac3n}+1}\to \frac32.
\]

\[
v_n=\sqrt{4n^2+n}-2n
=\frac{n}{\sqrt{4n^2+n}+2n}
=\frac1{\sqrt{4+\frac1n}+2}\to \frac14.
\]

\[
w_n=\sqrt{n^2+1}-\sqrt{n^2-1}
=\frac{2}{\sqrt{n^2+1}+\sqrt{n^2-1}}.
\]
Le dénominateur tend vers \(+\infty\), donc
\[
w_n\to0.
\]
}

\solution{Suites géométriques}{
Si \(|q|<1\), alors \(q^n\to0\).

Si \(q=1\), alors \(q^n=1\), donc la suite converge vers \(1\).

Si \(q=-1\), alors \(q^n=(-1)^n\), qui diverge.

Si \(q>1\), alors \(q^n\to+\infty\).

Si \(q<-1\), la suite n'a pas de limite réelle : les sous-suites paires et impaires ont des signes opposés et des valeurs absolues tendant vers \(+\infty\).
}

\solution{Passage à la limite dans les inégalités}{
Supposons par l'absurde que \(\ell>m\). Posons
\[
\eps=\frac{\ell-m}{3}>0.
\]
Pour \(n\) assez grand,
\[
u_n\geq \ell-\eps,\qquad v_n\leq m+\eps.
\]
Or
\[
\ell-\eps=\frac{2\ell+m}{3}>\frac{\ell+2m}{3}=m+\eps.
\]
Donc, pour \(n\) assez grand, \(u_n>v_n\), contradiction.

Exemple :
\[
u_n=0,\qquad v_n=\frac1n.
\]
Alors \(u_n<v_n\) pour tout \(n\geq1\), mais les deux suites convergent vers \(0\). On ne conserve donc que l'inégalité large.
}

\solution{Convergence implique bornitude}{
Puisque \(u_n\to\ell\), pour \(\eps=1\), il existe \(N\) tel que
\[
n\geq N\Rightarrow |u_n-\ell|\leq1.
\]
Alors
\[
|u_n|\leq |\ell|+1
\]
pour tout \(n\geq N\). Les termes \(u_0,\dots,u_{N-1}\) sont en nombre fini, donc leurs valeurs absolues admettent un maximum \(M_0\). En posant
\[
M=\max(M_0,|\ell|+1),
\]
on obtient
\[
\forall n\in\N,\quad |u_n|\leq M.
\]
Donc \((u_n)\) est bornée.

La réciproque est fausse : \(((-1)^n)\) est bornée mais divergente.
}

\solution{Monotonie par différence}{
Pour \(u_n=\frac n{n+1}\),
\[
u_{n+1}-u_n=\frac{n+1}{n+2}-\frac n{n+1}
=\frac1{(n+1)(n+2)}>0.
\]
Donc \((u_n)\) est croissante.

Pour \(v_n=n^2-3n\),
\[
v_{n+1}-v_n=(n+1)^2-3(n+1)-(n^2-3n)=2n-2.
\]
La suite décroît entre \(n=0\) et \(n=1\), puis est croissante pour \(n\geq1\).

Pour \(w_n=\frac1{n+1}\),
\[
w_{n+1}-w_n=\frac1{n+2}-\frac1{n+1}<0.
\]
Donc \((w_n)\) est décroissante.
}

\solution{Monotonie par quotient}{
On a
\[
\frac{u_{n+1}}{u_n}=\frac{2^{n+1}}{(n+1)!}\cdot \frac{n!}{2^n}=\frac2{n+1}.
\]
Donc
\[
\frac{u_{n+1}}{u_n}\geq1\Longleftrightarrow n+1\leq2\Longleftrightarrow n\leq1.
\]
La suite est croissante jusqu'au début, puis décroissante pour \(n\geq2\). Elle est positive. Comme elle devient décroissante positive, elle est bornée. Enfin, \(n!\) domine \(2^n\), donc
\[
\frac{2^n}{n!}\to0.
\]
Une justification élémentaire : pour \(n\geq4\),
\[
u_n=u_4\prod_{k=4}^{n-1}\frac2{k+1}\leq u_4\left(\frac25\right)^{n-4}\to0.
\]
}

\solution{Méthode \(\eps\)}{
On calcule :
\[
\left|\frac{2n+1}{3n-4}-\frac23\right|
=
\left|\frac{3(2n+1)-2(3n-4)}{3(3n-4)}\right|
=
\frac{11}{3(3n-4)}.
\]
Pour \(n\geq2\), \(3n-4>0\). On veut
\[
\frac{11}{3(3n-4)}\leq \eps.
\]
Il suffit de choisir \(N\) tel que
\[
3N-4\geq \frac{11}{3\eps}.
\]
Par exemple,
\[
N\geq \frac43+\frac{11}{9\eps}.
\]
Alors, pour tout \(n\geq N\),
\[
\left|\frac{2n+1}{3n-4}-\frac23\right|\leq\eps.
\]
Donc la limite est \(\frac23\).
}

\solution{Limites infinies}{
La définition est :
\[
u_n\to+\infty
\Longleftrightarrow
\forall A\in\R,\ \exists N\in\N,\ \forall n\geq N,\ u_n\geq A.
\]
Pour \(n\geq6\),
\[
n^2-3n=n(n-3)\geq \frac{n^2}{2}.
\]
Soit \(A\in\R\). Si \(A\leq0\), tout rang assez grand convient. Si \(A>0\), on choisit
\[
N\geq \max(6,\sqrt{2A}).
\]
Alors, pour \(n\geq N\),
\[
n^2-3n\geq \frac{n^2}{2}\geq A.
\]
Donc \(n^2-3n\to+\infty\).
}

\solution{Somme, produit, quotient}{
Par opérations sur les limites :
\[
a_n=3u_n-2v_n\to 3\cdot2-2(-3)=12.
\]
\[
b_n=u_nv_n\to 2\cdot(-3)=-6.
\]
Pour le quotient,
\[
v_n-1\to -4\neq0.
\]
Donc \(v_n-1\neq0\) à partir d'un certain rang, et
\[
c_n=\frac{u_n+1}{v_n-1}\to \frac{2+1}{-3-1}=-\frac34.
\]
}

\solution{Limite et valeur absolue}{
Si \(u_n\to0\), alors
\[
||u_n|-0|=|u_n|\to0.
\]
Donc \(|u_n|\to0\).

Réciproquement, si \(|u_n|\to0\), alors
\[
|u_n-0|=|u_n|\to0,
\]
donc \(u_n\to0\).

Pour \(\ell\neq0\), l'implication \(|u_n|\to|\ell|\Rightarrow u_n\to\ell\) est fausse. Exemple :
\[
u_n=-1,\qquad \ell=1.
\]
Alors \(|u_n|\to1=|\ell|\), mais \(u_n\to -1\neq1\).
}

\solution{Suite encadrée par deux suites convergentes}{
On a
\[
\frac{n}{n+1}\to1,
\qquad
\frac{n+2}{n+1}=1+\frac1{n+1}\to1.
\]
Par le théorème des gendarmes,
\[
u_n\to1.
\]
Un exemple est
\[
u_n=1,
\]
car
\[
\frac{n}{n+1}\leq1\leq\frac{n+2}{n+1}.
\]
}

\solution{Étude complète d'une suite explicite}{
On a
\[
u_n=\frac{n^2+1}{n^2+n+1}.
\]
Le numérateur et le dénominateur sont positifs, donc \(u_n>0\). De plus
\[
n^2+1<n^2+n+1
\]
pour \(n\geq1\), donc \(u_n<1\).

Monotonie :
\[
u_{n+1}-u_n=
\frac{(n^2+2n+2)(n^2+n+1)-(n^2+1)(n^2+3n+3)}
{(n^2+3n+3)(n^2+n+1)}.
\]
Le numérateur se simplifie en
\[
n^2+n-1.
\]
Pour \(n\geq1\), \(n^2+n-1>0\). Donc \((u_n)\) est croissante. Elle est majorée par \(1\), donc elle converge. Sa limite vaut
\[
\lim_{n\to\infty}\frac{1+\frac1{n^2}}{1+\frac1n+\frac1{n^2}}=1.
\]
}

\solution{Suite alternée simple}{
On a
\[
u_{2n}=1+\frac1{2n}\to1,
\]
et
\[
u_{2n+1}=-1+\frac1{2n+1}\to -1.
\]
Les deux suites extraites convergent vers deux limites distinctes. Donc \((u_n)\) ne converge pas.
}

\solution{Recouvrement pair-impair}{
Soit \(\eps>0\). Il existe \(N_1\) tel que
\[
n\geq N_1\Rightarrow |u_{2n}-\ell|\leq\eps,
\]
et \(N_2\) tel que
\[
n\geq N_2\Rightarrow |u_{2n+1}-\ell|\leq\eps.
\]
Posons \(N=2\max(N_1,N_2)+1\). Pour tout \(k\geq N\), soit \(k=2n\), soit \(k=2n+1\), avec \(n\geq\max(N_1,N_2)\). Donc
\[
|u_k-\ell|\leq\eps.
\]
Ainsi \(u_n\to\ell\).

Dans l'exemple,
\[
u_{2n}=\frac1{n+1}\to0,\qquad u_{2n+1}=\frac2{n+1}\to0.
\]
Donc \(u_n\to0\).
}

\solution{Bornes supérieure et inférieure}{
Pour tout \(n\geq1\),
\[
1-\frac1n\leq1.
\]
Donc \(1\) est un majorant. La suite
\[
a_n=1-\frac1n
\]
est une suite d'éléments de \(A\) et \(a_n\to1\). Ainsi, par caractérisation séquentielle,
\[
\sup A=1.
\]
Le plus petit élément de \(A\) est obtenu pour \(n=1\) :
\[
1-\frac11=0.
\]
Donc
\[
\inf A=0.
\]
}

\solution{Suite définie par une somme}{
On a
\[
\frac1{k(k+1)}=\frac1k-\frac1{k+1}.
\]
Donc
\[
u_n=\sum_{k=1}^n\left(\frac1k-\frac1{k+1}\right)
=1-\frac1{n+1}.
\]
La suite est croissante, car
\[
u_{n+1}-u_n=\frac1{n+1}-\frac1{n+2}>0.
\]
Elle converge vers
\[
1.
\]
}

\solution{Unicité de la limite}{
Supposons \(u_n\to\ell\) et \(u_n\to m\). Soit \(\eps>0\). Il existe \(N_1\) tel que, pour \(n\geq N_1\),
\[
|u_n-\ell|\leq\frac\eps2,
\]
et \(N_2\) tel que, pour \(n\geq N_2\),
\[
|u_n-m|\leq\frac\eps2.
\]
Pour \(n\geq\max(N_1,N_2)\),
\[
|\ell-m|\leq |\ell-u_n|+|u_n-m|\leq\eps.
\]
Donc
\[
\forall \eps>0,\quad |\ell-m|\leq\eps.
\]
Cela impose \(|\ell-m|=0\), donc \(\ell=m\).
}

\solution{Théorème des gendarmes}{
Soit \(\eps>0\). Comme \(a_n\to\ell\), il existe \(N_1\) tel que
\[
n\geq N_1\Rightarrow a_n\geq \ell-\eps.
\]
Comme \(b_n\to\ell\), il existe \(N_2\) tel que
\[
n\geq N_2\Rightarrow b_n\leq \ell+\eps.
\]
Enfin, il existe \(N_0\) tel que
\[
n\geq N_0\Rightarrow a_n\leq u_n\leq b_n.
\]
Pour \(n\geq\max(N_0,N_1,N_2)\),
\[
\ell-\eps\leq a_n\leq u_n\leq b_n\leq \ell+\eps.
\]
Donc
\[
|u_n-\ell|\leq\eps.
\]
Ainsi \(u_n\to\ell\).
}

\solution{Théorème de convergence monotone}{
L'ensemble
\[
A=\{u_n:n\in\N\}
\]
est non vide et majoré, donc il admet une borne supérieure \(\ell\). Par définition, pour tout \(n\),
\[
u_n\leq\ell.
\]
Soit \(\eps>0\). Puisque \(\ell=\sup A\), \(\ell-\eps\) n'est pas un majorant de \(A\). Donc il existe \(N\) tel que
\[
\ell-\eps<u_N\leq\ell.
\]
Comme \((u_n)\) est croissante, pour \(n\geq N\),
\[
u_N\leq u_n\leq\ell.
\]
Donc
\[
\ell-\eps<u_n\leq\ell,
\]
ce qui donne \(|u_n-\ell|\leq\eps\). Ainsi \(u_n\to\ell\).
}

\solution{Version divergente du théorème monotone}{
Soit \(A\in\R\). La suite n'étant pas majorée, il existe \(N\) tel que
\[
u_N>A.
\]
Comme \((u_n)\) est croissante, pour tout \(n\geq N\),
\[
u_n\geq u_N>A.
\]
Donc, pour tout \(A\in\R\), les termes de la suite sont supérieurs à \(A\) à partir d'un certain rang. Ainsi
\[
u_n\to+\infty.
\]
}

\solution{Suites adjacentes}{
Comme \(u_n\leq v_n\) et \((v_n)\) est décroissante, on a
\[
u_n\leq v_n\leq v_0.
\]
Donc \((u_n)\) est croissante et majorée ; elle converge vers une limite \(\ell\).

De même, comme \(u_0\leq u_n\leq v_n\), la suite \((v_n)\) est décroissante et minorée ; elle converge vers une limite \(m\).

Par opérations sur les limites,
\[
v_n-u_n\to m-\ell.
\]
Mais on sait que \(v_n-u_n\to0\). Par unicité de la limite,
\[
m-\ell=0.
\]
Donc \(\ell=m\).
}

\solution{Suites extraites d'une suite convergente}{
Comme \(\varphi\) est strictement croissante de \(\N\) dans \(\N\), on montre par récurrence que \(\varphi(n)\geq n\). C'est vrai pour \(n=0\). Si \(\varphi(n)\geq n\), alors
\[
\varphi(n+1)>\varphi(n).
\]
Comme ce sont des entiers,
\[
\varphi(n+1)\geq \varphi(n)+1\geq n+1.
\]

Soit \(\eps>0\). Comme \(u_n\to\ell\), il existe \(N\) tel que
\[
k\geq N\Rightarrow |u_k-\ell|\leq\eps.
\]
Si \(n\geq N\), alors \(\varphi(n)\geq n\geq N\), donc
\[
|u_{\varphi(n)}-\ell|\leq\eps.
\]
Ainsi \(u_{\varphi(n)}\to\ell\).
}

\solution{Critère de divergence par extraction}{
Si \((u_n)\) convergeait vers \(\ell\), alors toutes ses suites extraites convergeraient aussi vers \(\ell\). L'existence de deux suites extraites convergeant vers deux limites distinctes contredit donc la convergence de \((u_n)\).

Pour
\[
u_n=(-1)^n+\frac{1}{n+1},
\]
on a
\[
u_{2n}=1+\frac1{2n+1}\to1,
\]
et
\[
u_{2n+1}=-1+\frac1{2n+2}\to -1.
\]
Donc \((u_n)\) diverge.
}

\solution{Bolzano-Weierstrass : usage}{
Le théorème de Bolzano-Weierstrass affirme que toute suite réelle bornée possède une suite extraite convergente.

Comme \((u_n)\) est bornée, il existe donc une application strictement croissante \(\varphi:\N\to\N\) et un réel \(\ell\) tels que
\[
u_{\varphi(n)}\to\ell.
\]
Exemple de suite bornée divergente :
\[
u_n=(-1)^n.
\]
Elle est bornée par \(-1\) et \(1\), mais ses sous-suites paires et impaires convergent vers \(1\) et \(-1\).
}

\solution{Caractérisation séquentielle du supremum}{
Si \(s=\sup A\), alors \(s\) majore \(A\). Pour tout \(n\geq1\), \(s-\frac1n\) n'est pas un majorant de \(A\). Donc il existe \(a_n\in A\) tel que
\[
s-\frac1n<a_n\leq s.
\]
Alors
\[
0\leq s-a_n<\frac1n,
\]
donc, par encadrement,
\[
a_n\to s.
\]

Réciproquement, supposons que \(s\) majore \(A\), et qu'il existe \(a_n\in A\) tel que \(a_n\to s\). Soit \(M<s\). Posons \(\eps=s-M>0\). Pour \(n\) assez grand,
\[
|a_n-s|<\eps,
\]
donc
\[
a_n>s-\eps=M.
\]
Ainsi \(M\) ne majore pas \(A\). Tout réel strictement inférieur à \(s\) n'est pas majorant, et \(s\) est majorant. Donc
\[
s=\sup A.
\]
}

\solution{Comparaison asymptotique : définitions}{
On dit que
\[
u_n=o(v_n)
\]
si
\[
\frac{u_n}{v_n}\to0.
\]
On dit que
\[
u_n=O(v_n)
\]
si la suite \(\left(\frac{u_n}{v_n}\right)\) est bornée à partir d'un certain rang.
On dit que
\[
u_n\sim v_n
\]
si
\[
\frac{u_n}{v_n}\to1.
\]
Si \(u_n\sim v_n\), alors \(\frac{u_n}{v_n}\to1\), donc cette suite est bornée à partir d'un certain rang. Ainsi
\[
u_n=O(v_n).
\]
}

\solution{Une suite bornée dont les extraites convergent différemment}{
On calcule :
\[
u_{4n}=\sin(2n\pi)=0,
\]
\[
u_{4n+1}=\sin\left(2n\pi+\frac\pi2\right)=1,
\]
\[
u_{4n+2}=\sin((2n+1)\pi)=0,
\]
\[
u_{4n+3}=\sin\left(2n\pi+\frac{3\pi}{2}\right)=-1.
\]
La suite est bornée car \(|u_n|\leq1\). Elle diverge car deux suites extraites convergent vers \(1\) et \(-1\). L'ensemble des valeurs prises est
\[
\{-1,0,1\}.
\]
}

\solution{Suite récurrente affine}{
Le point fixe vérifie
\[
\ell=a\ell+b,
\]
donc
\[
\ell=\frac{b}{1-a}.
\]
Posons \(v_n=u_n-\ell\). Alors
\[
v_{n+1}=u_{n+1}-\ell=au_n+b-\ell.
\]
Or \(\ell=a\ell+b\), donc
\[
v_{n+1}=a(u_n-\ell)=av_n.
\]
Ainsi
\[
v_n=a^nv_0.
\]
Donc
\[
u_n=\ell+a^n(u_0-\ell).
\]
Comme \(|a|<1\), on a \(a^n\to0\), donc
\[
u_n\to \ell=\frac{b}{1-a}.
\]
}

\solution{Suite récurrente contractante simple}{
Si \(x\in[0,1]\), alors
\[
1\leq x+1\leq2,
\]
donc
\[
\frac12\leq f(x)\leq1.
\]
Ainsi \(f([0,1])\subset[0,1]\), et par récurrence \(u_n\in[0,1]\).

Ensuite,
\[
u_{n+1}-u_n=\frac{u_n+1}{2}-u_n=\frac{1-u_n}{2}\geq0.
\]
La suite est croissante et majorée par \(1\), donc elle converge. Soit \(\ell\) sa limite. Par passage à la limite,
\[
\ell=\frac{\ell+1}{2}.
\]
Donc
\[
2\ell=\ell+1,
\]
et
\[
\ell=1.
\]
}

\solution{Suite de Héron pour \(\sqrt2\)}{
Supposons \(u_n>\sqrt2\). Alors
\[
u_{n+1}-\sqrt2
=
\frac12\left(u_n+\frac2{u_n}\right)-\sqrt2
=
\frac{(u_n-\sqrt2)^2}{2u_n}>0.
\]
Donc \(u_{n+1}>\sqrt2\). Par récurrence,
\[
u_n>\sqrt2.
\]

De plus,
\[
u_{n+1}-u_n
=
\frac12\left(u_n+\frac2{u_n}\right)-u_n
=
\frac{2-u_n^2}{2u_n}.
\]
Comme \(u_n>\sqrt2\), on a \(2-u_n^2<0\). Donc
\[
u_{n+1}-u_n<0.
\]
La suite est décroissante et minorée par \(\sqrt2\), donc elle converge. Soit \(\ell\geq\sqrt2\) sa limite. Alors
\[
\ell=\frac12\left(\ell+\frac2\ell\right).
\]
Donc
\[
2\ell=\ell+\frac2\ell,
\]
puis
\[
\ell^2=2.
\]
Comme \(\ell>0\), on obtient
\[
\ell=\sqrt2.
\]
}

\solution{Deux suites couplées}{
On suppose \(0<u_n\leq v_n\). Alors
\[
u_{n+1}=\frac{2u_nv_n}{u_n+v_n}>0,\qquad v_{n+1}=\frac{u_n+v_n}{2}>0.
\]
L'inégalité entre moyenne harmonique et moyenne arithmétique donne
\[
u_{n+1}\leq v_{n+1}.
\]
Elle se vérifie directement :
\[
v_{n+1}-u_{n+1}
=
\frac{u_n+v_n}{2}-\frac{2u_nv_n}{u_n+v_n}
=
\frac{(v_n-u_n)^2}{2(u_n+v_n)}\geq0.
\]
Donc \(0<u_n\leq v_n\) pour tout \(n\).

On a
\[
u_{n+1}-u_n
=
\frac{2u_nv_n}{u_n+v_n}-u_n
=
\frac{u_n(v_n-u_n)}{u_n+v_n}\geq0.
\]
Donc \((u_n)\) est croissante.

De même,
\[
v_{n+1}-v_n
=
\frac{u_n+v_n}{2}-v_n
=
\frac{u_n-v_n}{2}\leq0.
\]
Donc \((v_n)\) est décroissante.

Enfin,
\[
v_{n+1}-u_{n+1}
=
\frac{(v_n-u_n)^2}{2(u_n+v_n)}.
\]
Comme \(u_n\geq u_0=a>0\), on a
\[
0\leq v_{n+1}-u_{n+1}
\leq \frac{(v_n-u_n)^2}{4a}.
\]
Les deux suites monotones bornées convergent vers \(\ell\) et \(m\), avec \(\ell\leq m\). En passant à la limite dans l'identité précédente,
\[
m-\ell=\frac{(m-\ell)^2}{2(\ell+m)}.
\]
Or \(\ell+m>0\). Si \(d=m-\ell\geq0\), alors
\[
d=\frac{d^2}{2(\ell+m)}.
\]
Comme \(d\leq b-a\) et l'itération contracte l'écart, on obtient nécessairement \(d=0\). Les suites sont adjacentes.
}

\solution{Suite définie par un maximum}{
Pour \(n\geq1\),
\[
u_{2n}=\max\left(\frac1{2n},\frac12\right)=\frac12.
\]
Et
\[
u_{2n+1}=\max\left(\frac1{2n+1},-\frac12\right)=\frac1{2n+1}.
\]
Donc
\[
u_{2n}\to\frac12,\qquad u_{2n+1}\to0.
\]
Les deux limites étant distinctes, \((u_n)\) diverge.
}

\solution{Suite et partie entière}{
Par définition,
\[
\lfloor \sqrt n\rfloor\leq \sqrt n,
\]
donc
\[
u_n\leq1.
\]
Comme \(n\geq1\), \(\lfloor \sqrt n\rfloor\geq1\), donc \(u_n>0\).

On sait que, pour tout réel \(x\),
\[
x-1<\lfloor x\rfloor\leq x.
\]
Avec \(x=\sqrt n\), on obtient
\[
\sqrt n-1<\lfloor\sqrt n\rfloor\leq\sqrt n.
\]
En divisant par \(\sqrt n>0\),
\[
1-\frac1{\sqrt n}<u_n\leq1.
\]
Les deux bornes tendent vers \(1\). Donc
\[
u_n\to1.
\]
}

\solution{Comparaison polynomiale}{
Le terme dominant de
\[
u_n=3n^4-2n^2+n
\]
est \(3n^4\). Donc
\[
u_n\sim3n^4,
\]
et \(u_n\to+\infty\).

Pour
\[
v_n=\frac{5n^3-n+7}{2n^2+1},
\]
on a
\[
v_n\sim \frac{5n^3}{2n^2}=\frac52 n.
\]
Donc \(v_n\to+\infty\).

Pour
\[
w_n=\frac{n^2+1}{n^4+n},
\]
on a
\[
w_n\sim \frac{n^2}{n^4}=\frac1{n^2}.
\]
Donc \(w_n\to0\).
}

\solution{Équivalents et racines}{
On écrit
\[
u_n=\sqrt{n^2+n+1}
=
n\sqrt{1+\frac1n+\frac1{n^2}}.
\]
Le facteur sous la racine tend vers \(1\), donc
\[
u_n\sim n.
\]

Ensuite,
\[
v_n=\sqrt{n+1}-\sqrt n
=
\frac1{\sqrt{n+1}+\sqrt n}.
\]
Comme le dénominateur est équivalent à \(2\sqrt n\),
\[
v_n\sim \frac1{2\sqrt n}.
\]

Enfin,
\[
w_n=\sqrt{n^2+2n}-n
=
\frac{2n}{\sqrt{n^2+2n}+n}
=
\frac2{\sqrt{1+\frac2n}+1}
\to1.
\]
Donc
\[
w_n\sim1.
\]
}

\solution{Contre-exemple sur les équivalents}{
Prenons
\[
u_n=n+1,\qquad v_n=n,
\]
alors \(u_n\sim v_n\). Prenons aussi
\[
a_n=-n,\qquad b_n=-n.
\]
Alors \(a_n\sim b_n\).

Mais
\[
u_n+a_n=(n+1)-n=1,
\]
tandis que
\[
v_n+b_n=n-n=0.
\]
On ne peut même pas former un équivalent de \(1\) avec \(0\). La compensation vient du fait que les termes principaux \(n\) et \(-n\) s'annulent dans les sommes.
}

\solution{Oral CCP -- Suite définie par \(u_{n+1}=u_n-u_n^2\)}{
Si \(0<u_n<1\), alors
\[
u_{n+1}=u_n(1-u_n).
\]
Comme \(0<u_n<1\), on a \(0<1-u_n<1\), donc
\[
0<u_{n+1}<u_n<1.
\]
Par récurrence, \(0<u_n<1\) pour tout \(n\). La suite est décroissante et minorée par \(0\), donc elle converge vers une limite \(\ell\geq0\).

En passant à la limite dans
\[
u_{n+1}=u_n-u_n^2,
\]
on obtient
\[
\ell=\ell-\ell^2,
\]
donc
\[
\ell=0.
\]

Ensuite,
\[
\frac1{u_{n+1}}-\frac1{u_n}
=
\frac1{u_n(1-u_n)}-\frac1{u_n}
=
\frac{1-(1-u_n)}{u_n(1-u_n)}
=
\frac1{1-u_n}.
\]
Comme \(u_n\to0\),
\[
\frac1{1-u_n}\to1.
\]
Ainsi
\[
\frac1{u_{n+1}}-\frac1{u_n}\to1.
\]
On en déduit par sommation de Cesàro que
\[
\frac1{u_n}\sim n,
\]
donc
\[
u_n\sim \frac1n.
\]
}

\solution{Oral Mines -- Suite logistique simple}{
Si \(0<u_n<1\), alors
\[
1<2-u_n<2,
\]
et
\[
u_{n+1}=u_n(2-u_n)>0.
\]
De plus,
\[
u_{n+1}<1
\Longleftrightarrow
u_n(2-u_n)<1
\Longleftrightarrow
1-2u_n+u_n^2>0
\Longleftrightarrow
(1-u_n)^2>0.
\]
Donc \(0<u_{n+1}<1\).

Monotonie :
\[
u_{n+1}-u_n=u_n(2-u_n)-u_n=u_n(1-u_n)>0.
\]
La suite est croissante et majorée par \(1\), donc elle converge. Soit \(\ell\) sa limite. Alors
\[
\ell=\ell(2-\ell),
\]
donc
\[
\ell-\ell(2-\ell)=0,
\]
c'est-à-dire
\[
\ell(\ell-1)=0.
\]
Comme \(u_n\) est croissante et \(u_0>0\), la limite est strictement positive ; donc
\[
\ell=1.
\]

Avec \(v_n=1-u_n\),
\[
v_{n+1}=1-u_{n+1}=1-u_n(2-u_n)=(1-u_n)^2=v_n^2.
\]
Ainsi
\[
v_n=v_0^{2^n}\to0,
\]
ce qui redonne \(u_n\to1\).
}

\solution{Oral Centrale -- Moyennes arithmético-géométriques}{
On utilise l'inégalité
\[
\sqrt{xy}\leq \frac{x+y}{2}
\]
pour \(x,y>0\). Si \(0<u_n\leq v_n\), alors
\[
0<u_{n+1}=\sqrt{u_nv_n}\leq \frac{u_n+v_n}{2}=v_{n+1}.
\]
Donc \(0<u_n\leq v_n\) pour tout \(n\).

Ensuite,
\[
u_{n+1}\geq u_n
\Longleftrightarrow
\sqrt{u_nv_n}\geq u_n
\Longleftrightarrow
v_n\geq u_n.
\]
Donc \((u_n)\) est croissante.

De même,
\[
v_{n+1}\leq v_n
\Longleftrightarrow
\frac{u_n+v_n}{2}\leq v_n
\Longleftrightarrow
u_n\leq v_n.
\]
Donc \((v_n)\) est décroissante.

Les deux suites sont monotones et bornées, donc convergent vers \(\ell\) et \(m\), avec \(\ell\leq m\). En passant à la limite dans
\[
u_{n+1}=\sqrt{u_nv_n},
\qquad
v_{n+1}=\frac{u_n+v_n}{2},
\]
on obtient
\[
\ell=\sqrt{\ell m},
\qquad
m=\frac{\ell+m}{2}.
\]
La seconde égalité donne directement \(m=\ell\). Les deux suites convergent donc vers la même limite.
}

\solution{Problème long -- Construction du nombre \(e\)}{
On admet ici les inégalités classiques permettant d'établir la monotonie :
\[
\left(1+\frac1n\right)^n
\leq
\left(1+\frac1{n+1}\right)^{n+1},
\]
donc \((u_n)\) est croissante, et
\[
\left(1+\frac1{n+1}\right)^{n+2}
\leq
\left(1+\frac1n\right)^{n+1},
\]
donc \((v_n)\) est décroissante.

On a immédiatement
\[
v_n=\left(1+\frac1n\right)u_n,
\]
donc
\[
u_n\leq v_n.
\]
De plus,
\[
v_n-u_n
=
u_n\left(1+\frac1n\right)-u_n
=
\frac{u_n}{n}.
\]
Comme \((u_n)\) est croissante et \(u_n\leq v_n\leq v_1=4\), elle est majorée. Donc
\[
0\leq v_n-u_n\leq \frac4n\to0.
\]
Ainsi \((u_n)\) et \((v_n)\) sont adjacentes. Elles convergent vers une même limite, appelée \(e\).
}

\solution{Problème long -- Récurrence et point fixe}{
Si \(x\in[0,2]\), alors
\[
2\leq 2+x\leq4,
\]
donc
\[
\sqrt2\leq f(x)\leq2.
\]
Ainsi \(f([0,2])\subset[0,2]\).

On étudie
\[
u_{n+1}-u_n=\sqrt{2+u_n}-u_n.
\]
Comme les deux termes sont positifs,
\[
\sqrt{2+u_n}\geq u_n
\Longleftrightarrow
2+u_n\geq u_n^2
\Longleftrightarrow
u_n^2-u_n-2\leq0
\Longleftrightarrow
(u_n-2)(u_n+1)\leq0.
\]
Sur \([0,2]\), on a \(u_n+1>0\) et \(u_n-2\leq0\). Donc
\[
u_{n+1}-u_n\geq0.
\]
La suite est croissante et majorée par \(2\), donc elle converge. Soit \(\ell\) sa limite. Alors
\[
\ell=\sqrt{2+\ell}.
\]
Donc
\[
\ell^2=\ell+2,
\]
c'est-à-dire
\[
(\ell-2)(\ell+1)=0.
\]
Comme \(\ell\in[0,2]\), on obtient
\[
\ell=2.
\]
}

\solution{Problème long -- Suites extraites et valeurs d'adhérence}{
Si
\[
u_{2n}\to\ell
\quad\text{et}\quad
u_{2n+1}\to\ell,
\]
alors, par le critère pair-impair, \(u_n\to\ell\).

Si les limites sont distinctes, la suite ne peut pas converger, car toute suite extraite d'une suite convergente converge vers la même limite.

Pour
\[
u_n=(-1)^n+\frac1n,
\]
on a
\[
u_{2n}=1+\frac1{2n}\to1,
\]
et
\[
u_{2n+1}=-1+\frac1{2n+1}\to-1.
\]
Donc la suite diverge.

Pour
\[
u_n=\frac{1+(-1)^n}{n+1},
\]
on a
\[
u_{2n}=\frac2{2n+1}\to0,
\]
et
\[
u_{2n+1}=0\to0.
\]
Donc
\[
u_n\to0.
\]
}

\solution{Problème long -- Une suite implicite par encadrement}{
Pour \(k\in\{1,\dots,n\}\), on a
\[
n+1\leq n+k\leq2n.
\]
Donc
\[
\frac1{2n}\leq \frac1{n+k}\leq \frac1{n+1}.
\]
En sommant,
\[
\frac n{2n}\leq u_n\leq \frac n{n+1}.
\]
Autrement dit,
\[
\frac12\leq u_n\leq \frac n{n+1}.
\]
Cet encadrement donne seulement
\[
\frac12\leq \liminf u_n\leq \limsup u_n\leq1.
\]
Il ne suffit pas à déterminer la limite. Une méthode plus fine consiste à reconnaître une somme de Riemann :
\[
u_n=\sum_{k=1}^n \frac1{n+k}
=
\frac1n\sum_{k=1}^n \frac1{1+k/n}.
\]
Hors du cadre strict des suites élémentaires, cette somme converge vers
\[
\int_0^1 \frac{dx}{1+x}=\ln2.
\]
}

\solution{Oral -- Suite complexe}{
Si \(|z|<1\), alors
\[
|z^n|=|z|^n\to0,
\]
donc \(z^n\to0\).

Si \(z=1\), alors \(z^n=1\), donc la suite converge vers \(1\).

Si \(|z|>1\), alors \(|z^n|=|z|^n\to+\infty\), donc la suite ne converge pas dans \(\C\).

Si \(|z|=1\), la convergence n'est pas automatique. Par exemple, pour \(z=-1\),
\[
z^n=(-1)^n
\]
diverge.
}

\solution{Oral -- Suite définie par \(u_{n+1}=\sqrt{u_n}\)}{
Par récurrence, \(u_n>0\) pour tout \(n\).

Si \(u_0=1\), alors \(u_n=1\) pour tout \(n\).

Si \(u_0>1\), alors
\[
\sqrt{u_n}<u_n,
\]
donc la suite est décroissante. Elle est minorée par \(1\), car si \(u_n>1\), alors \(\sqrt{u_n}>1\). Elle converge.

Si \(0<u_0<1\), alors
\[
\sqrt{u_n}>u_n,
\]
donc la suite est croissante. Elle est majorée par \(1\), car \(\sqrt{u_n}<1\). Elle converge.

Dans tous les cas, soit \(\ell\) la limite. Elle vérifie
\[
\ell=\sqrt{\ell}.
\]
Donc
\[
\ell^2=\ell,
\]
d'où
\[
\ell\in\{0,1\}.
\]
Comme \(u_n>0\) et que, dans les cas non constants, la suite est bornée du bon côté par \(1\), la limite est
\[
\ell=1.
\]
}

\solution{Oral -- Une suite rationnelle}{
On sait que
\[
1+2+\cdots+n=\frac{n(n+1)}2.
\]
Donc
\[
u_n=\frac{n(n+1)}{2n^2}=\frac{n+1}{2n}=\frac12+\frac1{2n}.
\]
Ainsi
\[
u_n\to\frac12.
\]
De plus,
\[
u_{n+1}-u_n
=
\left(\frac12+\frac1{2(n+1)}\right)
-
\left(\frac12+\frac1{2n}\right)
=
-\frac1{2n(n+1)}<0.
\]
La suite est donc strictement décroissante.
}

\end{document}